\documentclass{article}
\usepackage[utf8]{inputenc}
\usepackage{amsmath}
\usepackage{amsfonts}
\usepackage{amssymb}
\usepackage{graphicx}
\usepackage{hyperref}
\usepackage{booktabs}
\usepackage{longtable}
\usepackage{listings}
\usepackage{color}
\usepackage[margin=1in]{geometry}
\definecolor{lightgray}{gray}{0.95}
\lstset{
    backgroundcolor=\color{lightgray},
    basicstyle=\ttfamily\small,
    breaklines=true,
    frame=single,
}

\begin{document}

\section{Output Capabilities \& Typst Extensibility}

\subsection{1. Document Headers \& Configuration}


Headers, footers, and page numbers are intrinsically supported via YAML frontmatter blocks. For example, this document has a custom footer string.

\subsection{2. Advanced Typography \& Tables}


Because the intermediate format is Typst, typography is mathematically precise. Let's look at a standard Markdown table:

\begin{longtable}{lcrc}
\toprule
Subsystem & Underlying Tech & Role & Status\midrule
\textbf{Parser} & \texttt{pulldown-cmark} & Emits syntax events & ✅ \\
\textbf{Translator} & \texttt{marksmen-core} & Maps AST to Typst & ✅ \\
\textbf{Compiler} & \texttt{typst} & Typesets elements & ✅ \\
\textbf{Exporter} & \texttt{typst-pdf} & Emits binary format & ✅ \\
\bottomrule
\end{longtable}

\subsubsection{2.1 CSS \& HTML Handling}


Typst is an analytical typesetting engine, \textbf{not a web browser}. It cannot execute CSS or complex DOM layouts. However, the translator enforces semantic mapping for structural inline HTML subsets:

\begin{itemize}
\item \textbf{Underline}: This text is underlined via \texttt{\textless{}u\textgreater{}}
\item \textbf{Subscripts}: H2O via \texttt{\textless{}sub\textgreater{}}
\item \textbf{Superscripts}: E = mc2 via \texttt{\textless{}sup\textgreater{}}
\item \textbf{Colors}: This text is colored via \texttt{\textless{}span style="color: ..."\textgreater{}}
\item \textbf{Line breaks}: Forced HTML break  is translated correctly.
\end{itemize}


\textit{Note: Arbitrary \texttt{\textless{}div style="padding: 2px"\textgreater{}} blocks are ignored, preserving architectural invariants against empirical DOM approximations.}

\subsection{3. Mathematical Typesetting}


True native math (without KaTeX SVG hacks): \\

\[

\mathbf{J} = \nabla \times \mathbf{H} = \sigma \mathbf{E} + \frac{\partial \mathbf{D}}{\partial t}

\]



Inline equations like $\lim_{x \to \infty} f(x)$ blend seamlessly with the baseline.

\subsection{4. Linked SVG Figures}


Instead of relying on web-based JS to render graphics, Typst handles \texttt{SVG} vector embeddings natively without rasterization:


\begin{figure}[h]
\centering
\includegraphics[width=\textwidth]{./diagram1.svg}
Architecture Diagram\end{figure}


\subsection{5. Native Diagram Rendering (\texttt{marksmen-mermaid})}

\begin{quote}

\textbf{Pure-Rust Mermaid Translation:} \\
Executing standard \texttt{mermaid.js} typically requires evaluating an entire web runtime (V8/Browser) with empirical DOM measurement APIs, which violates our offline-first pure-Rust correctness guarantees.


\textit{Instead, \texttt{marksmen} relies on a mathematically-verified \texttt{marksmen-mermaid} layout engine implementing the Sugiyama Framework. It parses this raw Mermaid block offline and renders it deterministically into Typst spatial vectors!}

\end{quote}

\begin{lstlisting}[language=mermaid]
graph TD
    A[Markdown AST] --\textgreater{} B\{Marksmen Core\}
    B --\textgreater{} C[Typst Vectors]
    C --\textgreater{} D(PDF Output)
    C --\textgreater{} E(HTML5 Output)
\end{lstlisting}

\subsection{Appendix: Mathematical Justifications}


Because \texttt{marksmen} leverages Typst, we can seamlessly write $e^(i pi) + 1 = 0$ in the middle of text, and be confident that it renders beautifully on a completely separated page.

\end{document}
